\begin{DefinitionBox}{Parameterdarstellung einer Ebene}{analytic_geometry_parametric_plane}
\noindent
Eine Ebene $E$ mit Aufpunktvektor $\overrightarrow{OP}$ und \emph{2
linear unabhängigen Richtungsvektoren} $\vec a,\vec b\in\R^n$ besitzt die
Parameterdarstellung
\par\smallskip
\noindent
\begin{minipage}[t]{0.54\linewidth}
  \vspace{0pt}
  \centering
    \CachedTikzFigure{content/resources/chapters/analytic_geometry/tikz/parametric_plane_equation_annotation.tex}
\end{minipage}\hfill
\begin{minipage}[t]{0.42\linewidth}
  \vspace{0pt}
  \centering
  \CachedTikzFigure{content/resources/chapters/analytic_geometry/tikz/plane_parameter_representation.tex}
\end{minipage}
\par\smallskip
Die Parameter $\lambda$ und $\mu$ steuern unabhängig die Bewegung in den
beiden Richtungen der Ebene.
Nur linear unabhängige Richtungsvektoren $\vec a$ und $\vec b$ spannen
eine Ebene auf.
\end{DefinitionBox}

Wie bei Geraden kann der Aufpunkt durch jeden anderen Punkt der Ebene ersetzt
werden. Auch jedes andere Paar linear unabhängiger Richtungsvektoren in der
Ebene beschreibt dieselbe Ebene.
